\documentclass[10pt,twocolumn,letterpaper]{article}

\usepackage[utf8]{inputenc}
\usepackage[margin=0.75in]{geometry}
\usepackage{amsmath,amssymb,amsfonts}
\usepackage{booktabs}
\usepackage{graphicx}
\usepackage{microtype}
\usepackage{hyperref}
\usepackage{cite}
\usepackage{listings}
\usepackage{xcolor}

\hypersetup{
    colorlinks=true,
    linkcolor=blue,
    citecolor=blue,
    urlcolor=blue
}

\lstset{
    basicstyle=\ttfamily\footnotesize,
    breaklines=true,
    keywordstyle=\color{blue}\bfseries,
    stringstyle=\color{red!70!black},
    commentstyle=\color{gray},
    frame=single,
    tabsize=2
}

\title{\textbf{Do Unit Tests Guarantee Invariant Correctness?\\An Empirical Evaluation of Autonomous Program Repair on Distributed Systems Invariants}}

\author{
    \textbf{Tehya Lenka} \\
    Independent Research Project \\
    \texttt{lenkatehya@gmail.com} \\
    \small Artifact Repository: \url{https://github.com/tehyalenka/retrace-repair} \\
    \small Release Tag: \texttt{v1.9.0-research-release}
}

\date{November 2026}

\begin{document}

\maketitle

\begin{abstract}
Autonomous program repair (APR) systems powered by Large Language Models (LLMs) predominantly rely on unit test suites as their primary oracle for patch synthesis and validation. However, whether passing visible unit tests guarantees the preservation of complex software invariants remains an open empirical question. In this paper, we introduce \textbf{RETRACE}, an autonomous repair evaluation framework featuring physical dual-suite oracle isolation, disposable workspace replication, and cryptographic evidence sealing. We conduct an empirical study across \textbf{30 enterprise distributed systems tasks} spanning six architectural clusters: transactional durability, concurrency primitives, fault tolerance sagas, multi-tenant isolation, traffic resource pools, and coordination consensus. 

Across 180 controlled matrix evaluations evaluating two distinct model families (Groq \texttt{openai/gpt-oss-120b} and Alibaba \texttt{qwen/qwen3.8-27b}), we benchmark three repair paradigms: (1) stochastic \textbf{Resampling}, (2) standard iterative \textbf{Failure Feedback}, and (3) disposable \textbf{Counterfactual Replay}. Our central finding is that \textbf{public-test success substantially overestimates held-out invariant correctness across both model families}. While iterative feedback achieved an $83.3\%$--$86.7\%$ public pass rate (25--26/30 tasks), held-out invariant success was only $23.3\%$--$43.3\%$ (7--13/30 tasks). In $40.0\%$ of tasks on the 120B model and $63.3\%$ of tasks on the 27B model, LLM-generated patches successfully silenced visible test failures while silently breaking hidden architectural invariants. While Counterfactual Replay achieved the highest nominal success on the 120B model ($46.7\%$), paired exact McNemar tests indicate that the difference is not statistically significant ($p = 1.000$), while incurring an $80.5\%$ token cost premium. On the smaller 27B model, iterative feedback induced hyper-fixation on public traceback assertions, underperforming single-shot resampling ($23.3\%$ vs. $30.0\%$). We conclude with an empirical analysis of when speculative candidate exploration avoids local minima versus where localized feedback is computationally optimal.
\end{abstract}

\section{Introduction}

Automated Program Repair (APR) has experienced a paradigm shift from constraint-solving and genetic programming toward generative foundation models~\cite{fan2023automated,legoues2012genprog}. Modern coding agents typically operate via an iterative feedback loop: a model proposes code edits, a test harness executes the modified repository, and if execution fails, stack traces, compiler errors, and assertion failures are appended to the conversational context for subsequent refinement turns~\cite{xia2023keep,shinn2023reflexion}.

Despite impressive pass rates on synthetic code benchmarks~\cite{jimenez2024swebench}, a fundamental vulnerability persists: \textbf{oracle overfitting}~\cite{smith2015cure}. In production distributed systems, visible developer tests rarely constitute an exhaustive formal specification. A generated patch may satisfy isolated test assertions while violating broader systemic invariants, such as two-phase commit abort completeness, reader-writer lock fairness, lease expiry fencing, or multi-tenant boundary isolation.

When repair agents are evaluated exclusively against the test suites visible during prompt construction, the evaluation cannot distinguish between:
\begin{enumerate}
    \item \textbf{Generalizable Invariant Repair:} The patch restores the underlying behavioral contract of the system.
    \item \textbf{Oracle Exploitation:} The patch alters control flow specifically to silence failing assertions, introducing latent regressions.
\end{enumerate}

To investigate this phenomenon, we present \textbf{RETRACE}, a framework enforcing physical dual-suite isolation and cryptographic evidence sealing. We formulate three research questions:
\begin{itemize}
    \item \textbf{RQ1 (The Invariant Gap):} To what extent does passing visible unit tests predict success against held-out architectural invariant suites in distributed systems tasks?
    \item \textbf{RQ2 (Comparative Effectiveness):} Does speculative exploration via disposable workspace replay outperform standard failure feedback in discovering invariant-preserving repairs?
    \item \textbf{RQ3 (Computational Cost \& Tradeoffs):} What is the inference budget tradeoff (in model calls and token consumption) across repair paradigms?
\end{itemize}

\section{System Architecture}

RETRACE is architected around three core design requirements: physical oracle quarantine, disposable counterfactual exploration, and tamper-evident cryptographic provenance.

\subsection{Physical Dual-Suite Isolation}
To eliminate oracle leakage, RETRACE enforces physical filesystem boundaries:
\begin{itemize}
    \item \textbf{Public Test Suite (\texttt{tests/test\_public.py}):} Injected into the workspace, serving as visible CI feedback.
    \item \textbf{Held-Out Invariant Suite (\texttt{heldout.py}):} Quarantined outside the agent workspace in an isolated harness directory.
    \item \textbf{Reference Solution (\texttt{solution/}):} Stored out-of-tree for offline mutation hardening.
\end{itemize}
Neither \texttt{heldout.py} nor reference implementations are ever indexed, parsed, or exposed in model context windows. A formal test assertion (\texttt{test\_fixture\_isolation\_no\_heldout\_or\_solution\_leakage}) validates this isolation property prior to every execution cycle.

\subsection{Disposable Workspace Replay}
Rather than iteratively mutating a single mutable workspace—which risks polluting intermediate state—Counterfactual Replay provisions isolated filesystem clones in \texttt{/tmp}. Speculative candidate patches are generated, executed in disposable clones, evaluated against visible tests, and ranked deterministically. Only the highest-ranked candidate is emitted as the final patch.

\subsection{Cryptographic Integrity Envelope}
Prior to execution, RETRACE generates a pre-run manifest snapshot recording SHA-256 digests of prompts, repositories, and test fixtures. Post-execution, the runner executes a fail-closed verification pipeline (\texttt{research\_gate.py}, \texttt{statistical\_analysis.py}) and seals all 65 output artifacts into an immutable SHA-256 evidence manifest (\texttt{evidence-manifest.json}).

\section{Experimental Methodology}

\subsection{Benchmark Design \& Architectural Clusters}
The benchmark consists of \textbf{30 enterprise-grade distributed systems tasks} structured across six architectural clusters:
\begin{enumerate}
    \item \textbf{Transactional Durability (6 tasks):} WAL frame replay, two-phase commit (2PC) crash recovery, event sourcing aggregate versioning, order atomicity, ledger exactly-once delivery, and atomic bank balance transfers.
    \item \textbf{Concurrency \& CRDTs (5 tasks):} Distributed lock leases, optimistic concurrency control (OCC) document stores, fair reader-writer locks, PN-counters, and idempotent event consumer windows.
    \item \textbf{Fault Tolerance \& Sagas (6 tasks):} Dead-letter retries, cache stampede mutexes, bulkhead thread pools, outbox publishing, saga compensation coordinators, and circuit breaker state machines.
    \item \textbf{Multi-Tenant Isolation (5 tasks):} Secret rotation vaults, tenant-scoped token buckets, tamper-evident Merkle audit logs, RBAC policy evaluators, and isolated tenant caches.
    \item \textbf{Traffic \& Resource Pools (5 tasks):} Sliding-window rate limiters, database connection pools, burst token buckets, adaptive concurrency limiters, and SQLite-backed rate limiters.
    \item \textbf{Coordination \& Consensus (3 tasks):} Consistent hash rings with virtual nodes, leader epoch elections, and replicated state machine commit logs.
\end{enumerate}

\subsection{Mutation-Hardening Protocol}
To verify oracle sensitivity, all tasks were hardened using an automated AST mutation engine (\texttt{hardening.py}) applying four mutation operators: \texttt{return\_none}, \texttt{flip\_comparison}, \texttt{delete\_statement}, and \texttt{off\_by\_one}. Across 142 generated mutants, 116 were killed, achieving an \textbf{$81.69\%$ mutation score} with zero subprocess timeouts.

\subsection{Experimental Arms}
Each task is evaluated across three conditions:
\begin{itemize}
    \item \textbf{Resample:} Stochastic baseline sampled at $T = 0.7$ (seed $42$).
    \item \textbf{Feedback:} Iterative repair (up to 2 turns) at $T = 0.0$ incorporating failure tracebacks from public tests.
    \item \textbf{Counterfactual Replay:} Generates candidate repairs, evaluates them in disposable sandboxes, and selects the top patch.
\end{itemize}

\subsection{Model Specifications}
Evaluations were conducted across two model families via Groq endpoints:
\begin{itemize}
    \item \textbf{OpenAI GPT-OSS-120B:} 120-billion parameter open-weights model.
    \item \textbf{Alibaba Qwen-27B (\texttt{qwen/qwen3.8-27b}):} 27-billion parameter cross-architecture comparator.
\end{itemize}

\section{Empirical Results}

\begin{table*}[t]
\centering
\caption{Aggregate 30-Task Benchmark Results across Model Families ($N = 30$ tasks, 180 total evaluations).}
\label{tab:aggregate}
\begin{tabular}{llcccccc}
\toprule
\textbf{Strategy Arm} & \textbf{Model} & \textbf{Public Pass} & \textbf{Held-Out Pass} & \textbf{Public-Heldout Gap} & \textbf{Live Calls} & \textbf{Total Tokens} & \textbf{Tokens / Success} \\
\midrule
\textbf{Resample} & GPT-OSS-120B & 73.3\% (22/30) & \textbf{43.3\% (13/30)} & 30.0\% & 31 & 90,248 & \textbf{6,942.2} \\
                  & Qwen-27B     & 66.7\% (20/30) & \textbf{30.0\% (9/30)}  & 36.7\% & 36 & 73,285 & \textbf{8,142.8} \\
\midrule
\textbf{Feedback} & GPT-OSS-120B & \textbf{83.3\% (25/30)} & 43.3\% (13/30) & 40.0\% & 32 & 123,488 & 9,499.1 \\
                  & Qwen-27B     & \textbf{86.7\% (26/30)} & 23.3\% (7/30)  & \textbf{63.4\%} & 30 & 79,944  & 11,420.6 \\
\midrule
\textbf{Replay}   & GPT-OSS-120B & \textbf{83.3\% (25/30)} & \textbf{46.7\% (14/30)} & 36.6\% & 57 & 222,880 & 15,920.0 \\
                  & Qwen-27B     & \textbf{86.7\% (26/30)} & 20.0\% (6/30)  & \textbf{66.7\%} & 54 & 155,111 & 25,851.8 \\
\bottomrule
\end{tabular}
\end{table*}

\subsection{RQ1: The Public-to-Held-Out Invariant Gap}
Table~\ref{tab:aggregate} reports aggregate performance. Across both model families, \textbf{public unit-test success substantially overestimated held-out invariant correctness}:
\begin{itemize}
    \item On GPT-OSS-120B, Feedback achieved an $83.3\%$ public pass rate, but only a $43.3\%$ held-out pass rate (a $40.0\%$ drop). In 12 of 30 tasks, patches silenced visible CI tests while violating hidden invariants.
    \item On Qwen-27B, the gap widened to \textbf{$63.4\%$}: Feedback achieved an $86.7\%$ public pass rate (26/30), but only $23.3\%$ held-out invariant correctness (7/30). In 19 of 30 tasks, public tests gave false confidence.
\end{itemize}

\subsection{RQ2: Statistical Significance (McNemar Paired Tests)}
To assess whether Replay provides a statistically significant improvement over Feedback, we conducted two-tailed exact McNemar paired tests ($N = 30$)~\cite{mcnemar1947note}:
\begin{itemize}
    \item \textbf{GPT-OSS-120B:} 4 Replay-only wins vs. 3 Feedback-only wins ($p = 1.000$).
    \item \textbf{Qwen-27B:} 2 Replay-only wins vs. 3 Feedback-only wins ($p = 1.000$).
\end{itemize}
Because discordant pairs are symmetrically distributed, the evidence \textbf{does not support the hypothesis that Replay is generally superior to Feedback}.

\subsection{RQ3: Inference Budget Economics}
Replay incurred an \textbf{$80.5\%$ token premium} over Feedback on GPT-OSS-120B ($222,880$ vs. $123,488$ tokens; $15,920$ vs. $9,499$ tokens per held-out success). On Qwen-27B, Replay required $25,851$ tokens per held-out success compared to $8,143$ for Resample. The marginal cost of speculative search is substantial and must be weighed against defect topology.

\subsection{Cluster-Stratified Findings}
\begin{itemize}
    \item \textbf{Fault Tolerance \& Sagas:} Exhibited the largest invariant gap ($100\%$ public pass vs. $16.7\%$ held-out pass under Feedback).
    \item \textbf{Multi-Tenant Isolation:} Feedback was most effective ($80.0\%$ pass rate), as localized key filtering is easily resolved from tracebacks.
    \item \textbf{Traffic Pools:} Feedback failed completely ($0.0\%$ on Qwen, $0.0\%$ on 120B), where concurrency bursts exposed race conditions.
\end{itemize}

\section{Qualitative Failure Analysis}

\subsection{When Replay Prevents Oracle Overfitting}
In \texttt{enterprise\_circuit\_breaker}, public tests verified that repeated errors trip the state machine from \texttt{CLOSED} to \texttt{OPEN}. Under Feedback, the model observed a timeout failure and inserted a reset clearing \texttt{failure\_count} unconditionally upon half-open probe requests. While this satisfied visible tests, it broke the held-out invariant requiring probation failures to immediately reopen the circuit. Replay evaluated alternative candidates in disposable sandboxes, discarding overfitted branches and preserving probation invariants.

\subsection{When Feedback is More Efficient}
In \texttt{enterprise\_wal\_replay}, the failure was an unhandled \texttt{EOFError} during log truncation. The public traceback pinpointed the exact line of failure. Feedback resolved the bug in a single follow-up turn (1,200 tokens) via an explicit exception guard. Replay explored multi-file framing changes, consuming 8,000 tokens without improvement.

\section{Threats to Validity}

\begin{itemize}
    \item \textbf{Oracle Completeness:} While mutation hardening achieved an $81.7\%$ kill rate, held-out test suites are not formal proofs.
    \item \textbf{Model Scope:} Evaluated models comprise 120B and 27B open architectures. Closed-source reasoning models (e.g., o3-mini) may alter baseline pass rates.
    \item \textbf{Synthetic Tasks:} Repositories average 200--800 lines of systems code, evaluating algorithmic concurrency rather than multi-package dependency graphs.
\end{itemize}

\section{Conclusion}

Our empirical findings demonstrate that in enterprise systems repair, \textbf{passing public unit tests is an unreliable indicator of invariant correctness}. Across 180 evaluations on two model families, agents satisfied visible unit tests in $83.3\%$--$86.7\%$ of tasks, but preserved held-out invariants in only $20.0\%$--$46.7\%$. Replay offers targeted defense on complex state machines at the cost of an $80.5\%$ token premium, but does not achieve statistically significant dominance over feedback ($p = 1.000$). 

\bibliographystyle{plain}
\bibliography{references}

\end{document}
